\section{Step 7 Recognize the outer identifiers}\label{pcm-wav-s07}
The supplied frame retains input base in \texttt{s0}, complete count in \texttt{s1} and descriptor destination in \texttt{s2}. Compare bytes zero through three with RIFF and eight through eleven with WAVE. Use byte loads so an unaligned input pointer remains valid. Route a mismatch to \texttt{parse\_invalid}. This step identifies the container family; it says nothing yet about the encoding inside its chunks.

Author region \texttt{p\_magic}. Use \texttt{s07-1} and stop at \texttt{after\_p\_magic}.

\paragraph{Observe} Both identifiers match. Caller metadata is still unchanged.

\paragraph{Check} Does WAVE alone establish sixteen-bit PCM?
\begin{itemize}
\item[A.] Yes; every WAV has the same sample encoding.
\item[B.] Yes, when the filename ends in wav.
\item[C.] No; the format chunk must establish that.
\end{itemize}
